\section{7. Arquitectura y diseño}\label{arquitectura-y-diseuxf1o}

\subsection{7.1 Estilo arquitectónico}\label{estilo-arquitectuxf3nico}

El sistema sigue un estilo \textbf{cliente-servidor con API REST}. El servidor es un \textbf{monolito
modular en capas} (ADR-001):

\begin{itemize}
\tightlist
\item \textbf{Rutas} (\texttt{app/rutas/}): traducen HTTP a llamadas de servicio, validan la entrada con
Pydantic y declaran qué rol puede usar cada operación.
\item \textbf{Servicios} (\texttt{app/servicios/}): toda la lógica y las reglas de negocio RN-xx. Es la
única capa que cambia el estado de una orden o de una cotización.
\item \textbf{Modelo} (\texttt{app/models.py}): entidades del dominio, iguales a las clases de 6.3 y a las
tablas de \texttt{db/schema.sql}.
\item \textbf{Datos}: PostgreSQL (ADR-003), que además refuerza las restricciones con CHECK, índices únicos
parciales y disparadores.
\end{itemize}

Los clientes son dos interfaces web que consumen la misma API: la del personal del taller (tableta o
laptop) y el portal del cliente, al que se entra con el enlace seguro (ADR-005).
\pendiente{las interfaces web (PWA) se construyen en los Sprints 4 y 5; hoy la API se opera desde su
documentación interactiva en /docs.}

{\small
\begin{longtable}{@{}L{0.14\linewidth}L{0.82\linewidth}@{}}
\toprule
\textbf{RNF} & \textbf{Decisión de arquitectura que lo atiende} \\
\midrule
\endhead
\bottomrule
\endlastfoot
RNF-04 & Una sola capa de servicios registra las autorizaciones: guarda fecha y hora UTC, IP, medio, monto
y una huella SHA-256 de lo aceptado, que se recalcula para detectar alteraciones (CPA-20). El historial de
estados y los diagnósticos son registros inmutables (disparador \texttt{fn\_registro\_inmutable}). \\
RNF-05 & Contraseñas con PBKDF2-SHA256; tokens del cliente y códigos 2FA guardados solo como hash; control
de acceso por rol declarado en cada ruta; el portal identifica la orden por el token y no por el folio
(OWASP A01, CPA-30); HTTPS lo provee el hospedaje (ADR-004). \\
RNF-07 & Toda la pila es de código abierto: Python, FastAPI, SQLAlchemy, PostgreSQL, pytest (ADR-002 y
ADR-003). \\
RNF-08 y RNF-09 & Las entidades «OfflineSync» llevan \texttt{uuid\_local} único, generado en el dispositivo:
reenviar un registro ya sincronizado no lo duplica (sincronización idempotente). La captura sin conexión
vivirá en la PWA con IndexedDB. \pendiente{endpoint de sincronización por lotes (Sprint 4).} \\
RNF-10 & Las operaciones del portal son seguras de reintentar: pedir un código nuevo expira el anterior y
un código solo se acepta una vez (CPA-15), así que un reintento automático no duplica autorizaciones. \\
RNF-01 & El portal del cliente es web y responsivo; no requiere instalar nada ni conceder permisos. \\
\end{longtable}
}

\subsection{7.2 Diagramas C4}\label{diagramas-c4}

\figurapendiente{7.1}{C4 nivel 1: contexto del sistema}{el sistema y sus usuarios: Cliente (portal web
con enlace seguro), Recepcionista, Jefe de Mecánicos, Gerente y Encargado de Almacén; sistemas externos:
servicio de correo electrónico (enlaces y códigos 2FA) y, en Fase 2, servicio de SMS.}{docs/07-arquitectura/c4/contexto.puml}

\figurapendiente{7.2}{C4 nivel 2: contenedores}{PWA del personal (navegador en tableta/laptop, con
IndexedDB para operar sin conexión), portal del cliente (navegador móvil), API REST FastAPI (Python 3.12),
base de datos PostgreSQL 16, almacenamiento de fotografías y servicio de correo; protocolos HTTPS/JSON
entre navegador y API y SQL (psycopg) entre API y base de datos.}{docs/07-arquitectura/c4/contenedores.puml}

\figurapendiente{7.3}{C4 nivel 3: componentes de la API}{componentes Rutas del personal
(\texttt{rutas/personal.py}), Rutas del taller (\texttt{rutas/taller.py}), Portal (\texttt{rutas/portal.py}),
Seguridad (\texttt{seguridad.py}), Servicios de órdenes, citas, cotizaciones y cobro
(\texttt{servicios/}), Notificaciones (\texttt{notificaciones.py}) y Modelo/ORM
(\texttt{models.py}, \texttt{db.py}), con sus dependencias.}{docs/07-arquitectura/c4/componentes-api.puml}

\subsection{7.3 Registros de decisiones de arquitectura
(ADR)}\label{registros-de-decisiones-de-arquitectura-adr}

{\small
\begin{longtable}{@{}L{0.2\linewidth}L{0.76\linewidth}@{}}
\toprule
\textbf{Campo} & \textbf{Contenido} \\
\midrule
\endhead
\bottomrule
\endlastfoot
ID y título & \textbf{ADR-001 Estilo arquitectónico: monolito modular en capas con API REST} \\
Fecha y estado & 29/10/2026 · Aceptado (reemplaza la propuesta de microservicios de la Propuesta de
Solución Tecnológica) \\
Contexto & Equipo de 5 personas con poca experiencia en FastAPI (R-07), 6 sprints, hospedaje gratuito y
reglas de negocio que cruzan varias entidades en una misma transacción (autorizar cambia cotización,
orden e historial). \\
Alternativas evaluadas & 1. Microservicios (citas, órdenes, cotizaciones, cobro). 2. Monolito MVC con
páginas generadas en el servidor. 3. Monolito modular en capas que expone una API REST. \\
Decisión & Opción 3. Los módulos de servicio están separados por área y solo dependen hacia abajo
(7.5), de modo que podrían separarse después si hiciera falta. \\
Consecuencias & \textbf{+} Una sola transacción por operación; un solo despliegue; pruebas de extremo a
extremo sencillas. \textbf{+} La API sirve igual a la PWA del personal y al portal del cliente.
\textbf{--} Escala como una sola unidad (suficiente para un taller de 4 bahías). \textbf{--} Hay que vigilar
que las rutas no contengan lógica de negocio. \\
RNF relacionados & RNF-04, RNF-07, RNF-10 \\
\midrule
ID y título & \textbf{ADR-002 Pila tecnológica: Python 3.12 + FastAPI + SQLAlchemy 2} \\
Fecha y estado & 29/10/2026 · Aceptado (la interfaz web queda por decidir en el Sprint 3) \\
Contexto & La propuesta al cliente comprometió Python con FastAPI y herramientas gratuitas. Se necesita
validación de datos, documentación automática de la API y un ORM que funcione con PostgreSQL y SQLite. \\
Alternativas evaluadas & 1. Django + Django REST Framework. 2. Node.js + Express. 3. FastAPI + SQLAlchemy +
Pydantic. \\
Decisión & Opción 3: validación declarativa, documentación OpenAPI en \texttt{/docs} sin trabajo extra y
pruebas con \texttt{TestClient}. \\
Consecuencias & \textbf{+} Curva de aprendizaje corta para un equipo que ya usa Python. \textbf{--} FastAPI
no trae panel de administración ni autenticación integrada: se implementó la propia con la biblioteca
estándar (\texttt{seguridad.py}). \\
RNF relacionados & RNF-07, RNF-05 \\
\midrule
ID y título & \textbf{ADR-003 Base de datos: PostgreSQL 16} \\
Fecha y estado & 29/10/2026 · Aceptado \\
Contexto & El modelo tiene restricciones fuertes: especialización disjunta, una bahía con una sola orden
activa, una cita confirmada por bahía y horario, historial inmutable, atributos derivados. \\
Alternativas evaluadas & 1. MySQL 8. 2. SQLite. 3. PostgreSQL 16. \\
Decisión & PostgreSQL en el ambiente de pruebas y producción: tipos ENUM, índices únicos parciales,
columnas generadas y disparadores PL/pgSQL. SQLite solo para desarrollo local y pruebas automatizadas;
por eso las reglas también se validan en la capa de servicios. \\
Consecuencias & \textbf{+} La base de datos protege las reglas aunque alguien la modifique directamente.
\textbf{--} Dos motores (SQLite y PostgreSQL): la CI carga el esquema real en PostgreSQL para detectar
diferencias. \\
RNF relacionados & RNF-04, RNF-07, RNF-09 \\
\midrule
ID y título & \textbf{ADR-004 Hospedaje del ambiente de pruebas} \\
Fecha y estado & 29/10/2026 · \textcolor{red}{Propuesto (decidir antes del Sprint 4)} \\
Contexto & Se necesita un ambiente de pruebas con HTTPS, PostgreSQL y almacenamiento de fotografías, sin
costo (RNF-07). \\
Alternativas evaluadas & 1. Render (servicio web gratuito) + PostgreSQL gestionado gratuito (Neon o
Supabase). 2. Fly.io. 3. Servidor dentro del taller. \\
Decisión & \pendiente{elegir con el equipo; la recomendación inicial es la opción 1 por su HTTPS
automático y despliegue desde GitHub.} \\
Consecuencias & Opción 1: \textbf{+} sin servidor que mantener; \textbf{--} el plan gratuito suspende el
servicio tras inactividad y limita el almacenamiento (R-06). \\
RNF relacionados & RNF-05, RNF-07 \\
\midrule
ID y título & \textbf{ADR-005 Almacenamiento de fotografías y acceso del cliente con enlace seguro} \\
Fecha y estado & 29/10/2026 · Aceptado (acceso del cliente); fotografías: aceptado para v0.1.0 \\
Contexto & El cliente debe consultar y autorizar sin cuenta ni contraseña, pero nadie debe poder ver la
orden de otro (OWASP A01). Cada recepción genera al menos 4 fotografías que no deben saturar el
hospedaje gratuito. \\
Alternativas evaluadas & \textbf{Acceso:} 1. Cuenta de cliente con contraseña. 2. Folio + teléfono.
3. Token aleatorio en el enlace + código 2FA para autorizar. \textbf{Fotos:} 1. Guardarlas en la base de
datos (BLOB). 2. Sistema de archivos del servidor. 3. Servicio de almacenamiento de objetos. \\
Decisión & \textbf{Acceso:} opción 3 --- token de 256 bits generado con \texttt{secrets}, guardado solo como
hash SHA-256, un enlace activo por orden, que expira al entregar y saldar (RN-01); autorizar exige además
el código de 6 dígitos enviado al correo (RN-02). \textbf{Fotos:} opción 2 en v0.1.0 (carpeta
\texttt{MEDIA\_DIR}, solo JPEG/PNG/WebP de hasta 5 MB, nombre aleatorio); la base guarda solo la ruta, así
que migrar a la opción 3 no cambia el modelo. \\
Consecuencias & \textbf{+} El folio no sirve para entrar (CPA-30); quien intercepte el enlace aún necesita
el correo para autorizar. \textbf{--} Si el cliente pierde el enlace, el taller debe reenviarlo. \textbf{--}
Los archivos locales se pierden si el hospedaje gratuito reinicia el disco: \pendiente{confirmar con
ADR-004 y, de ser necesario, pasar a almacenamiento de objetos.} \\
RNF relacionados & RNF-04, RNF-05, RNF-01 \\
\midrule
ID y título & \textbf{ADR-006 Reglas de negocio como tablas de decisión DMN con Camunda} \\
Fecha y estado & 02/11/2026 · \textcolor{red}{Propuesto (diseño teórico; no implementado)} \\
Contexto & Las reglas RN-xx viven hoy como código Python en \texttt{app/servicios/}. Varias son
\emph{decisiones} con umbrales que el cliente puede querer cambiar (10\,\% de re-cotización, 5 días de
vigencia, tipos de vehículo admitidos, crédito a flotillas, conversión del nivel de combustible de RN-14).
Cada cambio de umbral exige hoy modificar código, probar y desplegar. \\
Alternativas evaluadas & 1. Mantener las reglas en Python (situación actual). 2. Tablas de decisión
DMN 1.3 modeladas en Camunda Modeler y evaluadas por el motor de decisiones de Camunda (DMN Engine) a
través de su API REST. 3. Tablas de decisión DMN evaluadas dentro del mismo proceso Python con una
biblioteca DMN. \\
Decisión & Opción 2 como diseño objetivo (detallado en 7.9): las decisiones se modelan en DMN con Camunda
Modeler, se versionan en \texttt{docs/07-arquitectura/dmn/} y la capa de servicios las consulta mediante
un adaptador \texttt{servicios/decisiones.py}. Los procesos (BPMN) siguen en el código; Camunda solo se
usa para \emph{decisiones}. \\
Consecuencias & \textbf{+} Las reglas quedan como tablas legibles que el Gerente puede revisar y validar
sin leer código. \textbf{+} Cada evaluación devuelve la regla que se cumplió, lo que sirve como evidencia
de auditoría (RNF-04). \textbf{--} Agrega un servicio más que desplegar y mantener. \textbf{--} El
motor no está disponible sin conexión: el adaptador debe tener un respaldo local (7.9). \textbf{--}
Hay que confirmar la licencia de la edición de Camunda que se use contra RNF-07 \pendiente{registrarlo como riesgo en 2.4}. \\
RNF relacionados & RNF-04, RNF-07, RNF-08 \\
\end{longtable}
}

\subsection{7.4 Diagramas de artefactos y de
despliegue}\label{diagramas-de-artefactos-y-de-despliegue}

\figurapendiente{7.4}{Diagrama de artefactos}{el artefacto \texttt{backend/} (paquete Python con
\texttt{app/} y \texttt{requirements.txt}), \texttt{db/schema.sql} y, cuando exista, el paquete estático de
la PWA; con sus dependencias de manifestación hacia los componentes de 7.3.}{docs/07-arquitectura/despliegue/artefactos.puml}

\figurapendiente{7.5}{Diagrama de despliegue del ambiente de pruebas}{nodo «dispositivo» tableta/laptop
del taller y celular del cliente (navegador), nodo «servidor» del hospedaje con el proceso uvicorn y la
API, nodo «base de datos» PostgreSQL 16 y almacenamiento de fotografías; conexiones HTTPS y TCP/5432. Debe
reflejar el hospedaje que se elija en ADR-004.}{docs/07-arquitectura/despliegue/despliegue.puml}

\subsection{7.5 Diagrama de paquetes}\label{diagrama-de-paquetes}

Dependencias entre paquetes de \texttt{backend/app} (todas en un solo sentido, sin ciclos):
\texttt{rutas} $\rightarrow$ \texttt{servicios} $\rightarrow$ \texttt{models} $\rightarrow$ \texttt{db};
\texttt{rutas} $\rightarrow$ \texttt{seguridad}, \texttt{schemas}; dentro de \texttt{servicios}:
\texttt{cobro} $\rightarrow$ \texttt{cotizaciones} $\rightarrow$ \texttt{ordenes}, y \texttt{citas} es
independiente.

\figurapendiente{7.6}{Diagrama de paquetes del backend}{los paquetes \texttt{app.rutas},
\texttt{app.servicios} (ordenes, citas, cotizaciones, cobro), \texttt{app.models}, \texttt{app.schemas},
\texttt{app.seguridad}, \texttt{app.notificaciones} y \texttt{app.db}, con las dependencias «import» descritas
arriba.}{docs/07-arquitectura/paquetes.puml}

\subsection{7.6 Diagramas de clases de la
solución}\label{diagramas-de-clases-de-la-soluciuxf3n}

\figurapendiente{7.7}{Clases de la solución del flujo núcleo}{las clases tal como existen en el código:
routers \texttt{personal}, \texttt{taller} y \texttt{portal}; módulos de servicio \texttt{ordenes}
(\texttt{abrir\_orden}, \texttt{asignar}, \texttt{registrar\_diagnostico}, \texttt{cambiar\_estado}),
\texttt{cotizaciones} (\texttt{crear\_cotizacion}, \texttt{solicitar\_codigo}, \texttt{confirmar\_codigo},
\texttt{huella}), \texttt{cobro} (\texttt{registrar\_pago}, \texttt{entregar}); entidades
\texttt{OrdenServicio}, \texttt{RecepcionVehicular}, \texttt{Cotizacion}, \texttt{Autorizacion},
\texttt{TokenAcceso}; y los esquemas Pydantic de entrada y salida.}{docs/07-arquitectura/clases-solucion.puml}

\subsection{7.7 Diagramas de secuencia}\label{diagramas-de-secuencia}

\figurapendiente{7.8}{Secuencia: recepción del vehículo (escenario principal del flujo núcleo)}{Recepcionista
→ \texttt{POST /ordenes} → \texttt{ordenes.abrir\_orden} (valida vehículo, cita y capacidad) → crea
OrdenServicio, HistorialEstado y RecepcionVehicular → respuesta con folio y alerta de bahías; luego
\texttt{POST /ordenes/\{folio\}/fotografias} por cada posición.}{docs/07-arquitectura/secuencias/recepcion.puml}

\figurapendiente{7.9}{Secuencia: autorización remota con código 2FA}{Cliente → \texttt{POST
/portal/\{token\}/codigo} → \texttt{token\_valido} → \texttt{solicitar\_codigo} → notificaciones (correo);
Cliente → \texttt{POST /portal/\{token\}/autorizacion} → \texttt{confirmar\_codigo} (vigencia, código,
huella) → \texttt{cambiar\_estado(Autorizada)}.}{docs/07-arquitectura/secuencias/autorizacion-remota.puml}

\figurapendiente{7.10}{Secuencia: re-cotización mayor al 10\,\%}{Jefe de mecánicos crea cotización
versión 2 → el servicio compara con el total inicial → queda Pendiente y se envía enlace →
\texttt{terminar} es rechazado (RN-04) hasta que el cliente autoriza.}{docs/07-arquitectura/secuencias/recotizacion.puml}

\figurapendiente{7.11}{Secuencia: pago y entrega}{Recepcionista → \texttt{POST pagos} →
\texttt{registrar\_pago} (no excede saldo) → \texttt{POST entrega} → \texttt{entregar} (saldo o crédito) →
Cerrada, Garantía, bahía libre y token expirado.}{docs/07-arquitectura/secuencias/entrega.puml}

\subsection{7.8 Diagrama de
comunicación}\label{diagrama-de-comunicaciuxf3n}

\figurapendiente{7.12}{Comunicación: asignación de bahía}{objetos \texttt{:RutaTaller},
\texttt{:ServicioOrdenes}, \texttt{os:OrdenServicio}, \texttt{b:Bahia}, \texttt{m:Usuario} y
\texttt{:HistorialEstado}, con los mensajes numerados 1 asignar, 1.1 obtener bahía, 1.2 validar
responsable, 1.3 ocupar bahía, 1.4 cambiarEstado(Abierta) si estaba en espera.}{docs/07-arquitectura/comunicacion-asignacion.puml}

\subsection{7.9 Modelado de decisiones con DMN y
Camunda}\label{modelado-de-decisiones-con-dmn-y-camunda}

\pendiente{esta sección describe el diseño propuesto en ADR-006. Camunda todavía no forma parte del
código ni del despliegue; las reglas siguen implementadas en Python en \texttt{app/servicios/}.}

\paragraph{Qué es DMN y para qué se usa aquí.} DMN (\emph{Decision Model and Notation}, estándar del
OMG, versión 1.3) describe decisiones de negocio como \textbf{tablas de decisión}: columnas de entrada,
columnas de salida y una fila por regla, con una \emph{política de acierto} (\emph{hit policy}) que dice qué
hacer si se cumplen varias filas. Las expresiones de cada celda se escriben en FEEL (\emph{Friendly
Enough Expression Language}), por ejemplo \texttt{> 10}, \texttt{[13..37]} o
\texttt{"Flotilla"}. \textbf{Camunda} aporta dos piezas: Camunda Modeler (aplicación de escritorio
gratuita para dibujar las tablas y guardarlas como archivos \texttt{.dmn}) y el motor de decisiones que
recibe los datos de entrada y devuelve el resultado.

En este proyecto Camunda se usa \textbf{solo para decisiones (DMN)}, no para orquestar procesos (BPMN): el
flujo de la orden (estados de 6.5) sigue controlado por \texttt{servicios/ordenes.py}, y únicamente le
pregunta al motor «¿qué dice la regla?» en los puntos donde hoy hay un \texttt{if} con umbrales de negocio.

\paragraph{Tablas de decisión propuestas.}

{\small
\begin{longtable}{@{}L{0.2\linewidth}L{0.14\linewidth}L{0.22\linewidth}L{0.18\linewidth}L{0.17\linewidth}@{}}
\toprule
\textbf{Decisión (clave)} & \textbf{Regla} & \textbf{Entradas} & \textbf{Salidas} & \textbf{Hit policy} \\
\midrule
\endhead
\bottomrule
\endlastfoot
\texttt{requiere\_recotizacion} & RN-04 & total inicial, total nuevo & requiere autorización (sí/no),
motivo & UNIQUE \\
\texttt{vigencia\_cotizacion} & RN-11 & días desde la emisión & vigente (sí/no) & UNIQUE \\
\texttt{vehiculo\_admitido} & RN-07 & tipo de vehículo, peso (t) & admitido (sí/no), motivo & FIRST \\
\texttt{permite\_entrega} & RN-10 & saldo pendiente, tipo de cliente, crédito autorizado por el Gerente &
entregar (sí/no), motivo & FIRST \\
\texttt{nivel\_combustible} & RN-14 & método de medición, porcentaje OBD-II & nivel (Reserva, 1/4, 1/2,
3/4, Lleno o No verificable) & UNIQUE \\
\texttt{nivel\_contingencia} & RN-15 & hay luz, hay internet, horas de corte, batería (\%) & nivel de
operación (1, 2 o 3) & FIRST \\
\end{longtable}
}

Como ejemplo, la tabla \texttt{permite\_entrega} (RN-10) queda así; se lee de arriba hacia abajo y gana la
primera fila que se cumpla (FIRST):

{\small
\begin{longtable}{@{}L{0.05\linewidth}L{0.17\linewidth}L{0.17\linewidth}L{0.2\linewidth}L{0.1\linewidth}L{0.21\linewidth}@{}}
\toprule
\textbf{\#} & \textbf{saldo} & \textbf{tipoCliente} & \textbf{creditoAutorizado} & \textbf{entregar} &
\textbf{motivo} \\
\midrule
\endhead
\bottomrule
\endlastfoot
1 & \texttt{<= 0} & \texttt{-} & \texttt{-} & \texttt{true} & \texttt{"Saldo cubierto"} \\
2 & \texttt{> 0} & \texttt{"Flotilla"} & \texttt{true} & \texttt{true} & \texttt{"Crédito a flotilla"} \\
3 & \texttt{> 0} & \texttt{-} & \texttt{-} & \texttt{false} & \texttt{"Saldo pendiente (RN-10)"} \\
\end{longtable}
}

\paragraph{Cómo se integra con el backend.}

\begin{enumerate}
\tightlist
\item \textbf{Modelado y versionado.} Cada tabla se dibuja en Camunda Modeler y se guarda en
\texttt{docs/07-arquitectura/dmn/<clave>.dmn}, en git, con revisión por Pull Request como el resto del
código. El Gerente valida la tabla impresa antes de publicarla.
\item \textbf{Despliegue.} Los archivos \texttt{.dmn} se despliegan al motor de Camunda, que corre como un
contenedor aparte junto a la API (un nodo más en el diagrama de despliegue de 7.4).
\item \textbf{Evaluación.} Un nuevo módulo \texttt{app/servicios/decisiones.py} expone una función
\texttt{evaluar(clave, datos)} que llama a la API REST del motor con las variables de entrada y devuelve
las salidas. Los demás servicios solo dependen de esa función; por ejemplo, \texttt{cobro.entregar} llama
a \texttt{evaluar("permite\_entrega", \{...\})} en lugar de comparar el saldo por su cuenta. Así se respeta
la regla de 7.1: la lógica de negocio sigue fuera de las rutas.
\item \textbf{Respaldo sin conexión.} Si el motor no responde (corte de internet, RNF-08, o contingencia
de RN-15), el adaptador evalúa una copia local de las mismas tablas con una biblioteca DMN de Python, o, en
su defecto, la implementación actual en Python. La PWA sin conexión no consulta al motor: las decisiones
que dependen de él se vuelven a evaluar al sincronizar (RNF-09).
\item \textbf{Auditoría.} Cada evaluación guarda en el historial de la orden la clave de la decisión, su
versión y la fila que se cumplió, como evidencia de por qué se permitió o bloqueó una operación
(RNF-04).
\item \textbf{Pruebas.} Cada fila de cada tabla tiene su caso en \texttt{backend/tests/}
(\texttt{test\_dmn\_<clave>.py}); los casos CPA existentes de RN-04, RN-07, RN-10 y RN-11 deben seguir
pasando sin cambios, lo que prueba que la migración no alteró el comportamiento.
\end{enumerate}

\paragraph{Plan de adopción.} Se migraría una decisión a la vez, empezando por \texttt{requiere\_recotizacion}
(RN-04), que es la más simple y la que el cliente tiene más probabilidad de ajustar. Antes de iniciar hay
que confirmar con el equipo y el docente: (1) la edición y versión de Camunda y que su licencia cumpla
RNF-07 para el uso que se le dará; (2) dónde se hospeda el motor (ADR-004); y (3) si el beneficio justifica
un servicio adicional en un taller de 4 bahías o si conviene la opción 3 de ADR-006 (DMN evaluado dentro de
la misma API).

\figurapendiente{7.13}{Diagrama de requisitos de decisión (DRD)}{las seis decisiones de la tabla de 7.9
como rectángulos de decisión, sus datos de entrada (orden, cotización, cliente, vehículo, recepción, estado
de la infraestructura) como óvalos y la fuente de conocimiento «Reglas de negocio 4.7» enlazada a cada
decisión; exportado de Camunda Modeler.}{docs/07-arquitectura/dmn/drd.dmn}
